\hypertarget{chapter-02-page-165}{}
\section{分岔理论与非唯一性结果}
\label{sec:ch2-bifurcation-nonuniqueness}

定常 Navier--Stokes 方程解的唯一性只在 $\nu$ 充分大, 或给定力与速度边界值充分小时得到证明. 可以预期, 否则解不唯一; \MBUnresolvedReference{citation}{V.I. Yudovich [1, 2]}, \MBUnresolvedReference{citation}{P. Rabinowitz [2]} 和 \MBUnresolvedReference{citation}{W. Velte [1, 2]} 已证明这一点. 前两篇论文处理 Bénard 问题(Navier--Stokes 方程与热传导方程), 第三篇证明 Taylor 问题的解不唯一.

本节依照 \MBUnresolvedReference{citation}{W. Velte [2]} 证明 Taylor 问题解的非唯一性. \MBUnresolvedReference{subsection}{Subsection 4.1} 描述问题并给出预备结果; \MBUnresolvedReference{subsection}{Subsection 4.2} 回顾所需的拓扑度理论主要结果; \MBUnresolvedReference{subsection}{Subsection 4.3} 随后证明非唯一性定理.

\subsection{Taylor 问题. 预备结果}
\label{subsec:ch2-taylor-preliminaries}

\paragraph{方程.}
\label{subsec:ch2-taylor-equations}
Taylor 问题研究 $\mathbb R^3$ 中一个区域内黏性不可压缩液体的流动; 该区域由半径为 $r_1,r_2$ 的两个无限长圆柱围成, $r_2>r_1>0$, 两圆柱具有同一竖直轴. 内圆柱以角速度 $\alpha$ 转动, 外圆柱静止.

由于寻找轴对称解, 使用 $\mathbb R^3$ 中的柱坐标 $r,\theta,z$, 其中 $0z$ 轴为圆柱轴. 流体充满区域 $\Omega$:
\begin{equation}
r_1<r<r_2,
\qquad-\infty<z<+\infty.
\label{eq:ch2-taylor-domain}
\end{equation}
以 $\widetilde u,\widetilde v,\widetilde w$ 表示速度向量在柱坐标中的分量, $\widetilde p$ 表示压力. 轴对称流运动方程的无量纲形式为
\begin{equation}
\left\{
\begin{aligned}
-\frac1\lambda\left(A\widetilde u-\frac{\widetilde u}{r^2}\right)
&+\widetilde u\frac{\partial\widetilde u}{\partial z}
+\widetilde w\frac{\partial\widetilde u}{\partial z}
-\frac1r\widetilde v^2+\frac{\partial\widetilde p}{\partial r}=0,\\
-\frac1\lambda\left(A\widetilde v-\frac{\widetilde v}{r^2}\right)
&+\widetilde u\frac{\partial\widetilde v}{\partial r}
+\widetilde w\frac{\partial\widetilde v}{\partial z}
+\frac1r\widetilde u\widetilde v=0,\\
-\frac1\lambda A\widetilde w
&+\widetilde u\frac{\partial\widetilde w}{\partial z}
+\widetilde w\frac{\partial\widetilde w}{\partial z}
+\frac{\partial\widetilde p}{\partial z}=0.
\end{aligned}
\right.
\label{eq:ch2-taylor-motion-equations}
\end{equation}
\begin{equation}
\frac\partial{\partial r}(r\widetilde u)
+\frac\partial{\partial z}(r\widetilde w)=0,
\label{eq:ch2-taylor-incompressibility}
\end{equation}
其中
\begin{equation}
A=\frac{\partial^2}{\partial r^2}
+\frac1r\frac\partial{\partial r}
+\frac{\partial^2}{\partial z^2},
\label{eq:ch2-taylor-operator-a}
\end{equation}
而 $\lambda=Re$ 是 Reynolds 数:
\begin{equation}
\lambda=Re=\alpha r_1^2\nu^{-1}.
\label{eq:ch2-taylor-reynolds-number}
\end{equation}
圆柱侧面上的边界条件为
\begin{equation}
\widetilde u=\widetilde w=0,\quad\widetilde v=1\quad\text{当 }r=r_1;
\qquad
\widetilde u=\widetilde v=\widetilde w=0\quad\text{当 }r=r_2.
\label{eq:ch2-taylor-boundary-data}
\end{equation}
方程 \eqref{eq:ch2-taylor-motion-equations}, \eqref{eq:ch2-taylor-incompressibility}, \eqref{eq:ch2-taylor-boundary-data} 有一个简单的已知解, 记为 $u_0,v_0,w_0,p_0$:
\begin{equation}
u_0=w_0=0,
\qquad
v_0(r)=\frac1{r_2^2-r_1^2}\left(\frac{r_2^2}{r}-r\right),
\qquad
p_0(r)=v_0^2\log r+\mathrm{const}.
\label{eq:ch2-taylor-trivial-base-solution}
\end{equation}
于是可寻找形如
\[
\widetilde u=u_0+u,
\qquad\widetilde v=v_0+v,
\qquad\widetilde w=w_0+w,
\qquad\widetilde p=p_0+p
\]
的解.

\hypertarget{chapter-02-page-167}{}
对 $u,v,w,p$ 得到方程
\begin{equation}
\left\{
\begin{aligned}
-\frac1\lambda\left(Au-\frac u{r^2}\right)
&+\frac{\partial u}{\partial r}
+w\frac{\partial u}{\partial z}
-\frac1rv^2-\frac2rvv_0+\frac{\partial p}{\partial r}=0,\\
-\frac1\lambda\left(Av-\frac v{r^2}\right)
&+u\frac{\partial v}{\partial r}
+w\frac{\partial v}{\partial z}
+\frac1ruv+\frac1r(v_0+rv_0')u=0,\\
-\frac1\lambda Aw
&+u\frac{\partial w}{\partial r}
+w\frac{\partial w}{\partial z}
+\frac{\partial p}{\partial z}=0,
\end{aligned}
\right.
\label{eq:ch2-taylor-perturbation-equations}
\end{equation}
其中 $v_0'=dv_0/dr$,
\begin{equation}
\frac\partial{\partial r}(ru)+\frac\partial{\partial z}(rw)=0,
\label{eq:ch2-taylor-perturbation-incompressibility}
\end{equation}
并且由 \eqref{eq:ch2-taylor-boundary-data} 得边界条件
\begin{equation}
u=v=w=0\quad\text{当 }r=r_1\text{ 或 }r=r_2.
\label{eq:ch2-taylor-homogeneous-boundary-data}
\end{equation}
注意, 条件 \eqref{eq:ch2-taylor-homogeneous-boundary-data} 不足以给出适定问题, 正如方程 \eqref{eq:ch2-taylor-motion-equations}, \eqref{eq:ch2-taylor-incompressibility} 的情形下条件 \eqref{eq:ch2-taylor-boundary-data} 也不充分; 还应在 $z=\pm\infty$ 加上某些边界条件. 可以寻找在 $z=\pm\infty$ 消失的解, 与 $z$ 无关的解, 或关于 $z$ 周期的解. 第一种可能性恰好对应第 \ref{sec:ch2-existence-uniqueness} 节研究的问题类型; 第二种也是如此, 因为它等价于寻找二维解. 第三种问题($z$-周期解)与第 \ref{sec:ch2-existence-uniqueness} 节所研究的问题非常接近, 并且这正是实验中观察到的解类型.

我们寻找 \eqref{eq:ch2-taylor-perturbation-equations}--\eqref{eq:ch2-taylor-homogeneous-boundary-data} 关于 $z$ 以 $L$ 为周期的解. 回忆 $u=v=w=p=0$ 是平凡解, 我们希望在某些情形下证明非平凡解存在.

\paragraph{流函数.}
\label{subsec:ch2-taylor-stream-function}
由 \eqref{eq:ch2-taylor-perturbation-incompressibility}, 存在函数 $f=f(r,z)$ 使
\begin{equation}
\frac\partial{\partial z}(rf)=ru,
\qquad
\frac\partial{\partial r}(rf)=-rw.
\label{eq:ch2-taylor-stream-function}
\end{equation}
边界条件 \eqref{eq:ch2-taylor-homogeneous-boundary-data} 保证 $f$ 是单值函数.

只用因变量 $f,v$ 写出 \eqref{eq:ch2-taylor-perturbation-equations}. 对第一式关于 $z$ 求导, 对第三式关于 $r$ 求导, 再将所得方程相减, 压力消失. 注意
\[
\frac{\partial u}{\partial z}-\frac{\partial w}{\partial r}=\mathcal Lf,
\qquad
\frac\partial{\partial r}(Aw)=\mathcal L\left(\frac{\partial w}{\partial r}\right),
\]
其中
\[
\mathcal L=A-\frac1{r^2}
=\frac{\partial^2}{\partial z^2}+\frac1r\frac\partial{\partial r}
+\frac{\partial^2}{\partial r^2}-\frac1{r^2}.
\]
展开后得到
\begin{equation}
\mathcal L^2f+\lambda\left(a\frac{\partial v}{\partial z}+M(f,v)\right)=0,
\label{eq:ch2-taylor-stream-equation}
\end{equation}

\hypertarget{chapter-02-page-168}{}
其中
\[
\begin{aligned}
M(f,v)
&=-\frac\partial{\partial r}\left(\frac{\partial f}{\partial z}\mathcal Lf\right)
+\frac\partial{\partial z}\left\{\frac1r\left(\frac\partial{\partial r}(rf)\right)\mathcal Lf\right\}
+\frac1r\frac\partial{\partial z}(v^2),\\
a(r)&=\frac{2v_0}{r}
=\frac2{r_2^2-r_1^2}\left(\frac{r_2^2}{r^2}-1\right).
\end{aligned}
\]
第二个方程可写成
\begin{equation}
\mathcal Lv+\lambda\left(b\frac{\partial f}{\partial z}+N(f,v)\right)=0,
\label{eq:ch2-taylor-azimuthal-equation}
\end{equation}
其中
\[
N(f,v)=-\frac1r\frac\partial{\partial r}(rv)\frac{\partial f}{\partial z}
+\frac1r\frac\partial{\partial r}(rf)\frac{\partial v}{\partial z},
\qquad
b=-\left(\frac{v_0}{r}+v_0'\right)=\frac2{r_2^2-r_1^2}.
\]
边界条件变为
\begin{equation}
f=\frac{\partial f}{\partial r}=v=0
\quad\text{当 }r=r_1\text{ 或 }r=r_2.
\label{eq:ch2-taylor-stream-boundary-data}
\end{equation}
问题于是归结为寻找 \eqref{eq:ch2-taylor-stream-equation}--\eqref{eq:ch2-taylor-stream-boundary-data} 关于 $z$ 以 $L$ 为周期的非平凡解 $f,v$.

\paragraph{相应的泛函方程.}
\label{subsec:ch2-taylor-functional-equation}
由于 $f,v$ 关于 $z$ 以 $L$ 为周期, 只需考虑它们在开集
\[
O_L=\{(r,z)\mid r_1<r<r_2,-L/2<z<L/2\}
\]
上的限制. 以 $H^k(O;L)$ 表示如下函数空间: 将其中的函数以周期 $L$ 延拓到整个集合
\[
O=\{(r,z)\mid r_1<r<r_2,z\in\mathbb R\},
\]
所得延拓函数属于 $O$ 的任意有界子集上的 $H^k$.\footnote{除零函数外, 这样的函数绝不属于 $H^k(O)$.}

实际上, $H^k(O;L)$ 正是 $H^k(O_L)$ 中满足
\[
\frac{\partial^jv}{\partial z^j}(r,-L/2)
=\frac{\partial^jv}{\partial z^j}(r,L/2),
\qquad j=0,\ldots,k-1
\]
的函数 $v$ 所成子空间. 引入空间 $W=W_1\times W_2$ 和 $V=V_1\times V_2$:
\[
\begin{aligned}
W_1&=\{f\in H^2(O;L)\mid f=\partial f/\partial r=0
\text{ 当 }r=r_1,r_2\},\\
W_2&=\{v\in H^1(O;L)\mid v=0\text{ 当 }r=r_1,r_2\},\\
V_1&=H^3(O;L)\cap W_1,\qquad V_2=W_2.
\end{aligned}
\]
$W_1,W_2,V_1,V_2$ 分别赋予由 $H^2(O_L),H^1(O_L),H^3(O_L),H^2(O_L)$ 诱导的范数后都是 Hilbert 空间.

\hypertarget{chapter-02-page-169}{}
在 $V$ 上采用由 $H^3(O_L)\times H^2(O_L)$ 诱导的积范数, 但在 $W$ 上赋予另一内积:
\begin{equation}
\begin{aligned}
((\phi_1,\phi_2))_W
&=((f_1,f_2))_{W_1}+((v_1,v_2))_{W_2},\\
((f_1,f_2))_{W_1}
&=\int_{O_L}\mathcal Lf_1\,\mathcal Lf_2\,r\,dr\,dz,\\
((v_1,v_2))_{W_2}
&=\int_{O_L}\left(
\frac{\partial v_1}{\partial z}\frac{\partial v_2}{\partial z}
+\frac{\partial v_1}{\partial r}\frac{\partial v_2}{\partial r}
\right)r\,dr\,dz,
\end{aligned}
\label{eq:ch2-taylor-w-inner-product}
\end{equation}
其中 $\phi_i=\{f_i,v_i\}\in W$. $W_2$ 上相应范数显然等价于 $H^1(O_L)$ 诱导的范数; 由下面的引理, $W_1$ 上相应范数等价于 $H^2(O_L)$ 的范数. 因而 $W$ 关于 \eqref{eq:ch2-taylor-w-inner-product} 是 Hilbert 空间.

\MBSourceStatementNumber{lemma}{1}
\begin{Lemma}
\label{lem:ch2-taylor-w1-norm-equivalence}
范数 $\lVert f\rVert_{W_1}$ 与 $H^2(O_L)$ 在 $W_1$ 上诱导的范数等价.
\end{Lemma}

\begin{Proof}
容易看出, 若 $f,g\in W_1$, 则
\begin{equation}
\int_{O_L}\mathcal Lf\,g\,r\,dr\,dz
=-\int_{O_L}\left(
\frac{\partial f}{\partial r}\frac{\partial g}{\partial r}
+\frac{\partial f}{\partial z}\frac{\partial g}{\partial z}
\right)r\,dr\,dz.
\label{eq:ch2-taylor-integration-by-parts}
\end{equation}
所以若 $f\in W_1$ 且 $\mathcal Lf=0$, 则 $\int_{O_L}\mathcal Lf\,f\,r\,dr\,dz=0$, 由 \eqref{eq:ch2-taylor-integration-by-parts} 得 $f=0$. 这表明 $\lVert f\rVert_{W_1}$ 是 $W_1$ 上的范数. 该范数等价于 $|\mathcal Lf|_{L^2(O_L)}$, 而由二阶椭圆算子 $-\mathcal L$ 的正则性定理, 后一范数等价于 $H^2(O_L)$ 的范数; 例如见 \MBUnresolvedReference{citation}{Agmon--Douglis--Nirenberg [1]}.
\end{Proof}

\paragraph{算子 $T$.}
在以下两个引理之后定义算子 $T$.

\MBSourceStatementNumber{lemma}{2}
\begin{Lemma}
\label{lem:ch2-taylor-elliptic-solvability}
对给定 $g\in L^2(O_L)$, 存在唯一的 $u\in W_2$(相应地 $v\in W_1$)使
\begin{equation}
\mathcal Lu=g
\label{eq:ch2-taylor-second-order-problem}
\end{equation}
(相应地
\begin{equation}
\mathcal L^2v=g).
\label{eq:ch2-taylor-fourth-order-problem}
\end{equation}
此外 $u\in H^2(O_L)$, $v\in H^4(O_L)$.
\end{Lemma}

\begin{Proof}
由 \eqref{eq:ch2-taylor-w-inner-product}, 这些问题分别等价于:

- 求 $u\in W_2$ 使
\[
-\int_{O_L}\left(
\frac{\partial u}{\partial r}\frac{\partial u_1}{\partial r}
+\frac{\partial u}{\partial z}\frac{\partial u_1}{\partial z}
\right)r\,dr\,dz
=\int_{O_L}gu_1r\,dr\,dz,
\qquad\forall u_1\in W_2.
\]

- 求 $v\in W_1$ 使
\[
\int_{O_L}\mathcal Lv\,\mathcal Lv_1\,r\,dr\,dz
=\int_{O_L}gv_1r\,dr\,dz,
\qquad\forall v_1\in W_1.
\]

\hypertarget{chapter-02-page-170}{}
这些变分问题解的存在唯一性由引理 \ref{lem:ch2-taylor-w1-norm-equivalence} 和投影定理 \ref{thm:ch1-abstract-projection} 得到. 再由椭圆算子 $\mathcal L,\mathcal L^2$ 的正则性定理, $u\in H^2(O_L)$, $v\in H^4(O_L)$.
\end{Proof}

\MBSourceStatementNumber{lemma}{3}
\begin{Lemma}
\label{lem:ch2-taylor-nonlinear-l2}
对给定 $\{f,v\}\in V$, 函数
\[
a\frac{\partial v}{\partial z}+M(f,v),
\qquad
b\frac{\partial f}{\partial z}+N(f,v)
\]
属于 $L^2(O_L)$.
\end{Lemma}

\begin{Proof}
只需证明 $M(f,v),N(f,v)\in L^2(O_L)$. Sobolev 嵌入定理特别给出
\[
H^2(O_L)\subset C^0(\overline O_L),
\qquad
H^1(O_L)\subset L^4(O_L).
\]
$M(f,v)$ 的首两项为
\[
-\frac{\partial^2f}{\partial r\partial z}\mathcal Lf
-\frac{\partial f}{\partial z}\frac\partial{\partial r}(\mathcal Lf).
\]
若 $\{f,v\}\in V$, 则 $f\in H^3(O_L)$, $\partial^2f/(\partial r\partial z)$ 和 $\mathcal Lf$ 属于 $H^1(O_L)$, 从而属于 $L^4(O_L)$, 二者之积属于 $L^2(O_L)$. 又 $\partial f/\partial z\in H^2(O_L)\subset C^0(\overline O_L)$, $\partial(\mathcal Lf)/\partial r\in L^2(O_L)$, 因而它们的乘积属于 $L^2(O_L)$. $M,N$ 中其余项的证明类似.
\end{Proof}

现在引入算子 $T$. 由引理 \ref{lem:ch2-taylor-elliptic-solvability}, \ref{lem:ch2-taylor-nonlinear-l2}, 对给定 $\{f,v\}\in V$, 存在唯一的解偶 $\{f',v'\}\in V\cap(H^4(O_L)\times H^2(O_L))$ 满足
\begin{equation}
\left\{
\begin{aligned}
\mathcal L^2f'&=-\left(a\frac{\partial v}{\partial z}+M(f,v)\right),\\
\mathcal Lv'&=-\left(b\frac{\partial f}{\partial z}+N(f,v)\right),\\
\{f',v'\}&\in W\quad(\text{或 }V).
\end{aligned}
\right.
\label{eq:ch2-taylor-operator-definition-system}
\end{equation}

\MBSourceStatementNumber{definition}{1}
\begin{Definition}
\label{def:ch2-taylor-operator-t}
以 $T$ 表示由 $\{f,v\}\mapsto\{f',v'\}$ 定义的从 $V$ 到自身的映射.
\end{Definition}

该算子与原问题的关系如下.

\MBSourceStatementNumber{proposition}{1}
\begin{Proposition}
\label{prop:ch2-taylor-functional-equivalence}
下列两个问题等价:

- 求 $\{f,v\}\in V$ 满足 \eqref{eq:ch2-taylor-stream-equation}, \eqref{eq:ch2-taylor-azimuthal-equation}, \eqref{eq:ch2-taylor-stream-boundary-data}.

- 求 $\phi=\{f,v\}\in V$ 使
\begin{equation}
\phi=\lambda T\phi.
\label{eq:ch2-taylor-functional-equation}
\end{equation}
\end{Proposition}

证明显然. 此后以泛函形式 \eqref{eq:ch2-taylor-functional-equation} 研究原问题. 下一段研究算子 $T$.

\hypertarget{chapter-02-page-171}{}
\paragraph{算子 $T$ 的性质.}
\label{subsec:ch2-taylor-operator-properties}

\MBSourceStatementNumber{lemma}{4}
\begin{Lemma}
\label{lem:ch2-taylor-operator-compact}
$T$ 是 $V$ 中的紧算子.
\end{Lemma}

\begin{Proof}
设 $\phi_n=\{f_n,v_n\}$ 是 $V$ 中的有界序列, $\phi_n'=\{f_n',v_n'\}=T\phi_n$. 引理 \ref{lem:ch2-taylor-nonlinear-l2} 的证明更精确地表明
\[
a\frac{\partial v_n}{\partial z}+M(f_n,v_n),
\qquad
b\frac{\partial f_n}{\partial z}+N(f_n,v_n)
\]
在 $L^2(O_L)$ 中有界. 因而 $\mathcal L^2f_n',\mathcal Lv_n'$ 在 $L^2(O_L)$ 中有界, $f_n',v_n'$ 分别在 $H^4(O_L),H^2(O_L)$ 中有界, 从而分别在 $H^3(O_L)$(或 $V_1$)和 $H^1(O_L)$(或 $V_2$)中相对紧.
\end{Proof}

考虑如下线性算子 $B$:
\begin{equation}
\begin{gathered}
\phi=\{f,v\}\longmapsto\phi''=\{f'',v''\}=B\phi,\\
\mathcal L^2f''=-a\frac{\partial v}{\partial z},
\qquad
\mathcal Lv''=-\frac{\partial f}{\partial z},
\qquad
\{f'',v''\}\in W.
\end{gathered}
\label{eq:ch2-taylor-linearized-operator}
\end{equation}
对给定 $\phi\in W$, 存在唯一的 $\{f'',v''\}\in W$ 满足 \eqref{eq:ch2-taylor-linearized-operator}. 因而 $B$ 是从 $W$ 到 $W\cap(H^4(O_L)\times H^2(O_L))$ 的连续线性映射.

\MBSourceStatementNumber{lemma}{5}
\begin{Lemma}
\label{lem:ch2-taylor-linearized-compact}
$B$ 在 $W$ 和 $V$ 中都是线性紧算子.
\end{Lemma}

\MBSourceStatementNumber{lemma}{6}
\begin{Lemma}
\label{lem:ch2-taylor-frechet-derivative}
算子 $B$ 是 $T$ 在点 $0$ 的 Fréchet 微分.
\end{Lemma}

\begin{Proof}
写成
\[
R\phi=\phi^*=T\phi-A\phi=\phi'-\phi'',
\qquad\phi\in V.
\]
于是
\begin{equation}
\mathcal L^2f^*=-M(f,v),
\qquad
\mathcal Lv^*=-N(f,v).
\label{eq:ch2-taylor-frechet-remainder-system}
\end{equation}
因为 $T0=0$, 必须证明
\begin{equation}
\frac{\lVert\phi^*\rVert_V}{\lVert\phi\rVert_V}\longrightarrow0
\quad\text{当 }\lVert\phi\rVert_V\to0.
\label{eq:ch2-taylor-frechet-remainder-limit}
\end{equation}
由 \eqref{eq:ch2-taylor-frechet-remainder-system} 和引理 \ref{lem:ch2-taylor-nonlinear-l2} 的方法,
\[
|M(f,v)|_{L^2(O_L)}\leq c_0\lVert\phi\rVert_V^2,
\qquad
|N(f,v)|_{L^2(O_L)}\leq c_1\lVert\phi\rVert_V^2.
\]
再由引理 \ref{lem:ch2-taylor-elliptic-solvability},
\[
\lVert f^*\rVert_{H^2(O_L)}\leq c_3\lVert\phi\rVert_V^2,
\qquad
\lVert v^*\rVert_{H^2(O_L)}\leq c_5\lVert\phi\rVert_V^2.
\]
特别地, $\lVert\phi^*\rVert_V\leq c_6\lVert\phi\rVert_V^2$, 从而 \eqref{eq:ch2-taylor-frechet-remainder-limit} 成立.
\end{Proof}

\hypertarget{chapter-02-page-172}{}
\paragraph{一个唯一性结果.}
\label{subsec:ch2-taylor-small-parameter-uniqueness}
在开始证明解的非唯一性之前, 先建立一个简单的唯一性结果($\lambda$ 小), 它正是把定理 \ref{thm:ch2-nonhomogeneous-uniqueness} 用于 Taylor 问题.

\MBSourceStatementNumber{proposition}{2}
\begin{Proposition}
\label{prop:ch2-taylor-small-parameter-uniqueness}
若 $\lambda$ 充分小, 即
\begin{equation}
0\leq\lambda<c(r_1,r_2),
\label{eq:ch2-taylor-small-parameter-condition}
\end{equation}
则问题 \eqref{eq:ch2-taylor-stream-equation}--\eqref{eq:ch2-taylor-stream-boundary-data}(或 \eqref{eq:ch2-taylor-functional-equation})在 $V$ 中除平凡解外没有其他解.
\end{Proposition}

\begin{Proof}
设 $\phi=\{f,v\}$ 是 \eqref{eq:ch2-taylor-stream-equation}--\eqref{eq:ch2-taylor-stream-boundary-data} 的解. 将 \eqref{eq:ch2-taylor-stream-equation} 乘 $f$, 将 \eqref{eq:ch2-taylor-azimuthal-equation} 乘 $v$, 再关于测度 $r\,dr\,dz$ 在 $O_L$ 上积分. 有
\[
\int_{O_L}[M(f,v)f-N(f,v)v]r\,dr\,dz=0.
\]
再用 \eqref{eq:ch2-taylor-integration-by-parts}, 得
\[
\int_{O_L}\left(|\mathcal Lf|^2+\left|\frac{\partial v}{\partial r}\right|^2
+\left|\frac{\partial v}{\partial z}\right|^2\right)r\,dr\,dz
=\lambda\int_{O_L}\left(-a\frac{\partial v}{\partial z}f
+b\frac{\partial f}{\partial z}v\right)r\,dr\,dz.
\]
分部积分可知右端等于
\[
\lambda\int_{O_L}(a+b)v\frac{\partial f}{\partial z}r\,dr\,dz.
\]
它不超过
\[
\lambda c(r_1,r_2)\int_{O_L}\left(|\mathcal Lf|^2
+\left|\frac{\partial v}{\partial r}\right|^2
+\left|\frac{\partial v}{\partial z}\right|^2\right)r\,dr\,dz,
\]
因为
\[
\sup_{r_1\leq r\leq r_2}|a+b|
=\frac2{r_2^2-r_1^2}\frac{r_2^2}{r_1^2},
\]
并且 Poincaré 不等式及引理 \ref{lem:ch2-taylor-w1-norm-equivalence} 分别给出
\[
|v|_{L^2(O_L)}^2\leq\mathrm{const}
\int_{O_L}\left(
\left|\frac{\partial v}{\partial r}\right|^2
+\left|\frac{\partial v}{\partial z}\right|^2\right)r\,dr\,dz,
\qquad
\left|\frac{\partial f}{\partial z}\right|_{L^2(O_L)}
\leq\mathrm{const}|\mathcal Lf|_{L^2(O_L)}.
\]
最终
\[
[1-\lambda c(r_1,r_2)]
\int_{O_L}\left(|\mathcal Lf|^2
+\left|\frac{\partial v}{\partial r}\right|^2
+\left|\frac{\partial v}{\partial z}\right|^2\right)r\,dr\,dz\leq0,
\]
当 \eqref{eq:ch2-taylor-small-parameter-condition} 成立时 $f=v=0$.\footnote{$c(r_1,r_2)$ 是依赖于 $r_1,r_2$ 的常数, 在命题证明中部分显式给出.}
\end{Proof}
\hypertarget{chapter-02-page-173}{}
\MBSourceStatementNumber{remark}{1}
\begin{Remark}
\label{rem:ch2-taylor-sharper-uniqueness}
稍加修改上述证明, 可证明当 $\lambda<\overline\lambda$ 时解唯一, 其中 $\overline\lambda$ 是算子 $(B+B^*)/2$ 的最小特征值, 该算子在 $W$ 中紧且自伴.
\end{Remark}

\MBSourceStatementNumber{remark}{2}
\begin{Remark}
\label{rem:ch2-taylor-large-parameter-goal}
我们的目标是证明在某些情形下, 当 $\lambda$ 充分大时解不唯一.
\end{Remark}

\subsection{$B$ 的一个特殊性质}
\label{subsec:ch2-taylor-special-property-b}

\paragraph{Fourier 级数展开.}
若 $f\in L^2(O_L)$, 则 $f$ 有 Fourier 级数展开
\begin{equation}
f=\sum_{n=0}^{\infty}
\bigl(f_n(r)\cos(n\sigma z)+\overline f_n(r)\sin(n\sigma z)\bigr),
\label{eq:ch2-taylor-fourier-l2-expansion}
\end{equation}
其中 $\sigma=2\pi/L$, $f_n,\overline f_n\in L^2(r_1,r_2)$. 此外
\[
\left\{\sum_{n=0}^{\infty}
\left(|f_n|_{L^2(r_1,r_2)}^2+|\overline f_n|_{L^2(r_1,r_2)}^2\right)\right\}^{1/2}
\]
有限, 并在 $L^2(O_L)$ 上定义一个与通常范数等价的范数.

$H^k(O;L)$ 中的函数($k\geq1$)也有 \eqref{eq:ch2-taylor-fourier-l2-expansion} 类型的展开, 其中 $f_n,\overline f_n\in H^k(r_1,r_2)$. 和
\begin{equation}
\left\{\sum_{n=0}^{\infty}\sum_{j=0}^{k}n^{2j}
\left(|f_n|_{H^{k-j}(r_1,r_2)}^2
+|\overline f_n|_{H^{k-j}(r_1,r_2)}^2\right)\right\}^{1/2}
\label{eq:ch2-taylor-fourier-sobolev-norm}
\end{equation}
有限, 并在 $H^k(O;L)$ 上定义一个与 $H^k(O_L)$ 诱导范数等价的范数.

考虑 $V$ 的子空间 $\overline V$, 其元素 $\{f,v\}$ 满足
\[
f(r,-z)=-f(r,z),
\qquad
v(r,-z)=v(r,z).
\]
这些函数有展开
\begin{equation}
f=\sum_{n=0}^{\infty}f_n\sin(n\sigma z),
\qquad
v=\sum_{n=0}^{\infty}v_n\cos(n\sigma z).
\label{eq:ch2-taylor-parity-fourier-expansion}
\end{equation}
前述结论容易限制到这些函数; $\overline V$ 是 $H^3(O;L)\times H^1(O;L)$ 的闭子空间.

若 $\phi=\{f,v\}\in\overline V$, 则 $\phi''=B\phi=\{f'',v''\}\in\overline V$, 且 $f'',v''$ 的 Fourier 系数由下列一维边值问题给出:
\begin{equation}
\left\{
\begin{aligned}
(\mathcal M-(n\sigma)^2)^2f_n''(r)&=an\sigma v_n(r),\\
-(\mathcal M-(n\sigma)^2)v_n''(r)&=bn\sigma f_n(r),\\
\mathcal M&=\frac{d^2}{dr^2}+\frac1r\frac d{dr}-\frac1{r^2},
\end{aligned}
\right.
\label{eq:ch2-taylor-fourier-boundary-system}
\end{equation}
边界条件为
\begin{equation}
f_n''(r)=\frac{df_n''}{dr}(r)=v_n(r)=0
\quad\text{当 }r=r_1\text{ 或 }r=r_2.
\label{eq:ch2-taylor-fourier-boundary-conditions}
\end{equation}

\hypertarget{chapter-02-page-174}{}
以 $B_n$ 表示线性映射 $\{f_n,v_n\}\mapsto\{f_n'',v_n''\}$; 它从 $L^2(r_1,r_2)^2$ 连续映入 $H^4(r_1,r_2)\times H^2(r_1,r_2)$.

\paragraph{算子 $B_n$ 的性质.}
下面的结果可见 \MBUnresolvedReference{citation}{S. Karlin [1]} 或 \MBUnresolvedReference{citation}{K. Kirchgässner [1]}.

\MBSourceStatementNumber{lemma}{7}
\begin{Lemma}
\label{lem:ch2-taylor-green-function-signs}
设
\[
\mathcal M=\alpha(r)\frac{d^2}{dr^2}
+\beta(r)\frac d{dr}+\gamma(r),
\]
其中 $\alpha,\beta,\gamma$ 在 $[-1,+1]$ 上四次连续可微, $\alpha>0$, $\gamma<0$. 则:

(i) 在边界条件 $G(\pm1,s)=0$ 下, $\mathcal M$ 的 Green 函数 $G(r,s)$ 对 $r,s\in(-1,+1)$ 为负;

(ii) 在边界条件 $H(\pm1,s)=\partial H/\partial z(\pm1,s)=0$ 下, $\mathcal M^2$ 的 Green 函数 $H(r,s)$ 对 $r,s\in(-1,+1)$ 为正.
\end{Lemma}

由此推出:

\MBSourceStatementNumber{lemma}{8}
\begin{Lemma}
\label{lem:ch2-taylor-shifted-green-positive}
在边界条件 \eqref{eq:ch2-taylor-fourier-boundary-conditions} 下, $(\mathcal M-k^2)^2$ 和 $-(\mathcal M-k^2)$ 在 $(r_1,r_2)$ 上的 Green 函数 $H(k;r,s),G(k;r,s)$ 在 $(r_1,r_2)^2$ 上为正.
\end{Lemma}

利用这些核, 可将 \eqref{eq:ch2-taylor-fourier-boundary-system}, \eqref{eq:ch2-taylor-fourier-boundary-conditions} 化为积分方程
\begin{equation}
f_n''(r)=n\sigma\int_{r_1}^{r_2}G(n\sigma;r,s)a(s)v_n(s)\,ds,
\label{eq:ch2-taylor-integral-f-double-prime}
\end{equation}
\begin{equation}
v_n''(r)=n\sigma\int_{r_1}^{r_2}H(n\sigma;r,s)bf_n(s)\,ds.
\label{eq:ch2-taylor-integral-v-double-prime}
\end{equation}
$B_n$ 的特征向量是满足 $B_n\{f_n,v_n\}=\lambda\{f_n,v_n\}$ 的函数偶. 因而 \eqref{eq:ch2-taylor-fourier-boundary-system}, \eqref{eq:ch2-taylor-fourier-boundary-conditions} 中 $f_n''=\lambda f_n$, $v_n''=\lambda v_n$. 这等价于
\begin{equation}
f_n(r)=\lambda n\sigma\int_{r_1}^{r_2}G(n\sigma;r,z)a(s)v_n(s)\,ds,
\label{eq:ch2-taylor-eigen-integral-f}
\end{equation}
\begin{equation}
v_n(r)=\lambda n\sigma\int_{r_1}^{r_2}H(n\sigma;r,s)bf_n(s)\,ds.
\label{eq:ch2-taylor-eigen-integral-v}
\end{equation}
消去 $v_n$ 还得到
\begin{equation}
f_n(r)=\mu\int_{r_1}^{r_2}K(n\sigma;r,s)f_n(s)\,ds,
\label{eq:ch2-taylor-reduced-eigenproblem}
\end{equation}
其中 $\mu=\lambda^2$, 且
\begin{equation}
K(k;r,s)=k^2\int_{r_1}^{r_2}G(k;r,t)H(k;t,s)a(t)b\,dt.
\label{eq:ch2-taylor-positive-kernel}
\end{equation}

\hypertarget{chapter-02-page-175}{}
下面给出算子 $B_n$ 特征值的一些性质. 它们基于 \MBUnresolvedReference{citation}{Witting [1]} 中的结果.\footnote{这是关于保持 Banach 空间中一个锥不变的线性紧算子的 \MBUnresolvedReference{citation}{Krein--Rutman [1]} 一般结果的特例. 这些结果是线性代数中正矩阵的 Perron--Frobenius 定理在无限维的推广; 例如见 \MBUnresolvedReference{citation}{R.S. Varga [1]}.}

\MBSourceStatementNumber{lemma}{9}
\begin{Lemma}
\label{lem:ch2-taylor-positive-kernel-eigenvalue}
设 $K(r,s)$ 是定义在正方形 $[r_1,r_2]^2$ 上的实连续函数, 且在其内部严格为正. 特征值问题
\begin{equation}
f(r)=\lambda\int_{r_1}^{r_2}K(r,s)f(s)\,ds,
\qquad r_1<r<r_2,
\label{eq:ch2-taylor-positive-kernel-eigenproblem}
\end{equation}
有一个解 $\lambda_1>0$, 对应的特征函数 $f\in C^0([r_1,r_2])$ 且在 $r_1<r<r_2$ 时 $f(r)>0$. \eqref{eq:ch2-taylor-positive-kernel-eigenproblem} 的任意其他特征解对应于满足 $|\lambda|>\lambda_1$ 的特征值 $\lambda$.
\end{Lemma}

\MBSourceStatementNumber{lemma}{10}
\begin{Lemma}
\label{lem:ch2-taylor-bn-principal-eigenvalue}
算子 $B_n$ 有特征值 $\lambda_n^1>0$, 对应的特征向量 $\{f_n^1,v_n^1\}$ 在 $r_1<r<r_2$ 时满足 $f_n^1(r)>0,v_n^1(r)>0$. $B_n$ 的任意其他特征值 $\lambda$ 满足 $|\lambda|>\lambda_n^1$.
\end{Lemma}

\begin{Proof}
由 \eqref{eq:ch2-taylor-reduced-eigenproblem}, $\{f_n,v_n\}$ 是 $B_n$ 的特征值为 $\lambda_n$ 的特征向量, 当且仅当 $f_n$ 是 \eqref{eq:ch2-taylor-reduced-eigenproblem} 的特征解且 $\mu=\lambda_n^2$. 引理 \ref{lem:ch2-taylor-positive-kernel-eigenvalue} 可用于 \eqref{eq:ch2-taylor-reduced-eigenproblem}, 除 $v_n^1$ 的正性外均得证. 后者由引理 \ref{lem:ch2-taylor-shifted-green-positive}, $f_n^1$ 的正性以及把 \eqref{eq:ch2-taylor-eigen-integral-v} 写成 $v_n=v_n^1,f_n=f_n^1$ 得到.
\end{Proof}

\MBSourceStatementNumber{lemma}{11}
\begin{Lemma}
\label{lem:ch2-taylor-bn-spectral-gap}
设 $\lambda_n$ 是 $B_n$ 的不同于 $\lambda_n^1$ 的特征值. 则
\begin{equation}
|\lambda_n|>\lambda_n^1
\geq\frac2{\max|a+b|}\,n^2\sigma^2.
\label{eq:ch2-taylor-bn-spectral-bound}
\end{equation}
\end{Lemma}

\begin{Proof}
第一式由引理 \ref{lem:ch2-taylor-bn-principal-eigenvalue} 得到. 为证明第二式, 有
\begin{equation}
\left\{
\begin{aligned}
(\mathcal M-(n\sigma)^2)^2f_n^1(r)&=\lambda_n^1an\sigma v_n^1(r),\\
-(\mathcal M-(n\sigma)^2)v_n^1(r)&=\lambda_n^1bn\sigma f_n^1(r).
\end{aligned}
\right.
\label{eq:ch2-taylor-principal-eigen-system}
\end{equation}
将第一式乘 $rf_n^1(r)$, 第二式乘 $rv_n^1(r)$, 再积分并分部积分, 得
\begin{equation}
J(f_n^1,v_n^1)
=\lambda_n^1n\sigma\int_{r_1}^{r_2}(a+b)f_n^1v_n^1r\,dr,
\label{eq:ch2-taylor-principal-energy-identity}
\end{equation}
其中
\[
\begin{aligned}
J(f,v)=\int_{r_1}^{r_2}\biggl[&(\mathcal Mf)^2
+2(n\sigma)^2\left(\left(\frac{df}{dr}\right)^2+\frac1{r^2}f^2\right)
+(n\sigma)^4f^2\\
&+\left(\frac{dv}{dr}\right)^2+\frac1{r^2}v^2+(n\sigma)^2v^2\biggr]r\,dr.
\end{aligned}
\]

\hypertarget{chapter-02-page-176}{}
式 \eqref{eq:ch2-taylor-principal-energy-identity} 的右端不超过
\[
\begin{aligned}
&\lambda_n^1n\sigma\max|a+b|
\int_{r_1}^{r_2}|f_n^1v_n^1|r\,dr\\
&\leq\lambda_n^1n\sigma\max\frac{|a+b|}{2}
\left\{n\sigma\int_{r_1}^{r_2}|f_n^1|^2r\,dr
+\frac1{n\sigma}\int_{r_1}^{r_2}|v_n^1|^2r\,dr\right\}\\
&\leq\frac{\lambda_n^1}{(n\sigma)^2}
\max\frac{|a+b|}{2}J(f_n^1,v_n^1).
\end{aligned}
\]
因为 $J(f_n^1,v_n^1)\ne0$, 可约去它, 得到 \eqref{eq:ch2-taylor-bn-spectral-bound}.
\end{Proof}

\paragraph{$B$ 的谱性质.}
首先注意, 若 $\{f,v\}$ 是 $B$ 的特征值为 $\lambda$ 的特征向量, 则其展开 \eqref{eq:ch2-taylor-parity-fourier-expansion} 对应的 $\{f_n,v_n\}$ 分别是 $B_n$ 的特征值为 $\lambda$ 的特征向量. 反过来由上一段导出 $B$ 的谱性质.

考虑引理 \ref{lem:ch2-taylor-bn-principal-eigenvalue} 给出的全部 $\lambda_n^1$, $n\geq1$. 由 \eqref{eq:ch2-taylor-bn-spectral-bound}, $\lambda_n^1\to\infty$, 因而 $\inf_{n\geq1}\lambda_n^1$ 有限且严格为正. 以 $m$ 表示满足 $\lambda_m^1=\inf_{p\geq1}\lambda_p^1$ 的最大整数 $n$. 可能对另一些 $n<m$ 也有 $\lambda_m^1=\lambda_n^1$. 下一个引理避开这种情形.

\MBSourceStatementNumber{lemma}{12}
\begin{Lemma}
\label{lem:ch2-taylor-simple-principal-eigenvalue}
可选择周期 $L$ 使
\[
\lambda_1^1>\lambda_n^1,
\qquad\forall n>1.
\]
在这种情形下, 记 $\lambda_1^1=\inf_{n\geq1}\lambda_n^1$ 为 $\lambda_1$; $\lambda_1$ 是 $B$ 在 $\overline V$ 中的简单特征值, $B$ 的任意其他特征值 $\lambda$ 满足 $|\lambda|>\lambda_1$. 对应于 $\lambda_1$ 的特征向量 $\phi_1=\{f^1,v^1\}$ 有 \eqref{eq:ch2-taylor-parity-fourier-expansion} 类型的展开, 其中对所有 $n>1$, $f_n^1=v_n^1=0$.
\end{Lemma}

\begin{Proof}
任取 $L$, 令 $\sigma=2\pi/L$, $m$ 如上定义. 置 $\sigma^*=m\sigma$, $L^*=L/m$. 对相应算子 $B$, $\lambda_1$ 是 $B_1$ 的特征值, 且对所有 $n>1$ 有 $\lambda_1<\lambda_n^1$. 其余性质显然.
\end{Proof}

本节余下部分假设 $L$ 的选择使引理 \ref{lem:ch2-taylor-simple-principal-eigenvalue} 成立.

最后一个结果涉及 $\lambda_1$ 作为 $B$ 的特征值的次数. 该次数是 $\operatorname{Ker}(I-\lambda_1B)^p$ 的维数, 当 $p$ 充分大时与 $p$ 无关.

\MBSourceStatementNumber{lemma}{13}
\begin{Lemma}
\label{lem:ch2-taylor-principal-degree-one}
在引理 \ref{lem:ch2-taylor-simple-principal-eigenvalue} 的条件下, $\lambda_1$ 是 $B$ 的次数为 $1$ 的特征值.
\end{Lemma}

\begin{Proof}
证明若 $p\geq2$ 且 $(I-\lambda_1B)^p\phi=0$, 则 $(I-\lambda_1B)\phi=0$. 于是对每个 $p$, $\operatorname{Ker}(I-\lambda_1B)^p=\operatorname{Ker}(I-\lambda_1B)$, 由引理 \ref{lem:ch2-taylor-simple-principal-eigenvalue}, 其维数为 $1$. 对 $p$ 归纳, 只需证明 $(I-\lambda_1B)^2\phi=0$ 蕴含 $(I-\lambda_1B)\phi=0$. 取 $\phi^0$ 使
\[
(I-\lambda_1B)^2\phi^0=0.
\]

\hypertarget{chapter-02-page-177}{}
反证. 假设 $(I-\lambda_1B)\phi^0$ 不等于 $0$. 于是这个向量在相差一个乘法常数的意义下等于先前的特征函数 $\phi^1$:
\begin{equation}
\phi^1=(I-\lambda_1B)\phi^0,
\qquad
\phi^1=\lambda_1B\phi^1.
\label{eq:ch2-taylor-generalized-eigenvector-relations}
\end{equation}
有
\[
\phi^0=\lambda_1B(\phi^0+\phi^1),
\]
这等价于
\begin{equation}
\mathcal L^2f^0=\lambda_1a\frac\partial{\partial z}(v^0+v^1),
\qquad
\mathcal Lv^0=\lambda_1b\frac\partial{\partial z}(f^0+f^1).
\label{eq:ch2-taylor-generalized-eigenvector-system}
\end{equation}
回忆
\[
f^1(r,z)=f_1^1(r)\sin(\sigma z),
\qquad
v^1(r,z)=v_1^1(r)\cos(\sigma z).
\]
考虑 $f^0,v^0$ 的 Fourier 级数:
\[
f^0=\sum_{n=1}^{\infty}f_n^0\sin(n\sigma z),
\qquad
v^0=\sum_{n=1}^{\infty}v_n^0\cos(n\sigma z).
\]
关系 \eqref{eq:ch2-taylor-generalized-eigenvector-system} 蕴含
\begin{equation}
\left\{
\begin{aligned}
(\mathcal M-\sigma^2)^2f_1^0&=\lambda_1a\sigma(v_1^0+v_1^1),\\
-(\mathcal M-\sigma^2)v_1^0&=\lambda_1b\sigma(f_1^0+f_1^1),
\end{aligned}
\right.
\label{eq:ch2-taylor-first-mode-generalized-system}
\end{equation}
而当 $n\geq2$ 时,
\[
(\mathcal M-\sigma^2)^2f_n^0=\lambda_1a\sigma n v_n^0,
\qquad
-(\mathcal M-\sigma^2)v_n^0=\lambda_1b\sigma n f_n^0.
\]
因为当 $n\geq2$ 时 $\lambda_1$ 不是 $B_n$ 的特征值, 故 $f_n^0=v_n^0=0$.

现在像处理 \eqref{eq:ch2-taylor-integral-f-double-prime}, \eqref{eq:ch2-taylor-integral-v-double-prime} 那样, 将 \eqref{eq:ch2-taylor-first-mode-generalized-system} 化为积分方程. 因为 $\{f_1^1,v_1^1\}$ 是 $B_1$ 的特征向量, 得
\[
f_1^0(r)=\lambda_1\int_{r_1}^{r_2}G(\sigma;r,s)a(s)\sigma v_1^0(s)\,ds+f_1^1(r),
\]
\[
v_1^0(r)=\lambda_1\int_{r_1}^{r_2}H(\sigma;r,s)b\sigma f_1^0(s)\,ds+v_1^1(r).
\]
消去 $v_1^0$, 并使用 \eqref{eq:ch2-taylor-positive-kernel} 中引入的核 $K$, 得
\begin{equation}
f_1^0(r)-\lambda_1^2\int_{r_1}^{r_2}K(\sigma;r,s)f_1^0(s)\,ds
=2f_1^1(r).
\label{eq:ch2-taylor-fredholm-contradiction-equation}
\end{equation}
方程 \eqref{eq:ch2-taylor-fredholm-contradiction-equation} 满足 Fredholm 二择一. 因此 $f_1^1$ 与伴随方程的特征函数 $g_1$ 正交:
\[
g_1(r)-\lambda_1^2\int_{r_1}^{r_2}K(\sigma;s,r)g_1(s)\,ds=0.
\]
由引理 \ref{lem:ch2-taylor-positive-kernel-eigenvalue}, $f_1^1,g_1$ 在 $(r_1,r_2)$ 上为正, 这与正交条件
\[
\int_{r_1}^{r_2}f_1^1(s)g_1(s)\,ds=0
\]
矛盾. 因此 $(I-\lambda_1B)\phi^0=0$, 证明完成.
\end{Proof}

\hypertarget{chapter-02-page-178}{}
\subsection{拓扑度理论的若干要素}
\label{subsec:ch2-topological-degree-elements}
下面回顾拓扑度理论的若干定义和性质. 关于证明与进一步结果, 读者可参见 \MBUnresolvedReference{citation}{J. Leray and J. Schauder [1]}, 或 \MBUnresolvedReference{citation}{M. A. Krasnoselskii [1]}, \MBUnresolvedReference{citation}{L. Nirenberg [1]}, \MBUnresolvedReference{citation}{P. Rabinowitz [4]} 的基础工作.

\paragraph{拓扑度.}
设 $T$ 是赋范空间 $V$ 中的紧算子, $S=I-T$(其中 $I$ 是 $V$ 中的恒等算子). 以 $\omega,\omega_i$ 表示 $V$ 的有界区域, 以 $\overline\omega,\partial\omega$ 表示 $\omega$ 的闭包和边界.

若 $\omega$ 是 $V$ 的有界区域, $v\in V$, 且
\[
v\notin S(\partial\omega),
\]
则可定义一个整数 $d(S,\omega,v)$, 称为 $S$ 在 $\omega$ 中于点 $v$ 的拓扑度.

拓扑度的主要性质如下:

(i) 若 $\omega=\omega_1\cup\omega_2$, $\omega_1\cap\omega_2=\varnothing$, 且 $v\notin S(\partial\omega_1)$, $v\notin S(\partial\omega_2)$, 则 $v\notin S(\partial\omega)$, 且
\[
d(S,\omega,v)=d(S,\omega_1,v)+d(S,\omega_2,v).
\]

(ii) 若 $d(S,\omega,v)\ne0$, 则 $v\in S(\omega)$, 也就是说方程
\[
(I-T)(u)=v
\]
在 $\omega$ 中至少有一个解.

(iii) 若 $S,\omega,v$ 连续变化, 且 $v$ 始终不属于 $S(\partial\omega)$, 则 $d(S,\omega,v)$ 保持不变.\footnote{$S$ 的连续变化定义如下: $S=S(\lambda)=I-T(\lambda)$, $\lambda\in\mathbb R$(或任意拓扑空间), 且映射 $\lambda\mapsto T(\lambda)\phi$ 关于 $\phi$(其中 $\phi\in V$)一致连续.}

\paragraph{指数.}
设 $u_0$ 是 $V$ 中一点, $v=Su_0$, 并假设方程 $Su=v$ 在 $u_0$ 的某个邻域中只有解 $u_0$.

在这种情形下, 对充分小的 $\epsilon$, 可以定义度 $d(S,\omega_\epsilon(u_0),v)$, 其中 $\omega_\epsilon(u_0)$ 是以 $u_0$ 为中心、半径为 $\epsilon$ 的开球. 由拓扑度的性质 (iii), 当 $\epsilon\to0$ 时该数与 $\epsilon$ 无关.

于是把 $S$ 在 $u_0$ 处的指数定义为充分小 $\epsilon$ 时的这个度:
\[
i(S,u_0)=d(S,\omega_\epsilon(u_0),v),
\qquad \epsilon<\epsilon_0.
\]
下面列出该指数的一些基本性质:

(i) 若方程 $Su=v$ 在有界区域 $\omega$ 中有有限多个解 $u_k$, 并且在 $\partial\omega$ 上无解, 则
\[
d(S,\omega,v)=\sum_k i(S,u_k).
\]

(ii) 恒等算子($T=0$)在任意点 $u_0$ 的指数为 $1$: $i(I,u_0)=1$.

(iii) 假设 $T$ 在点 $u_0$ 有 Fréchet 微分 $A$. 则 $A$ 与 $T$ 一样是紧算子. 若 $I-A$ 是一一的(即 $1$ 不是 $A$ 的特征值), 则 $u_0$ 是方程 $Su=Su_0$ 的孤立解, 且可定义指数 $i(S,u_0)$. 有
\[
i(S,u_0)=i(I-A,0)=i(I-A).
\]

\hypertarget{chapter-02-page-179}{}
(iv) 若 $A$ 是 $V$ 中的线性紧算子, 且 $I-A$ 是一一的, 则 $I-A$ 的指数为 $\pm1$.\footnote{当且仅当 $I-A$ 是一一的时, 才能定义 $I-A$ 的指数; 定义后, 它在每一点 $u_0$ 都相同.}

类似地, 可在不含 $A$ 的任何特征值的任意区间 $\lambda'\leq\lambda\leq\lambda''$ 上定义 $I-\lambda A$ 的指数; 该指数在这样的区间上为常数, 并等于 $\pm1$. 特别地, 在区间 $(0,\lambda_1)$ 上 $i(I-\lambda A)=1$, 其中 $\lambda_1$ 是 $A$ 的最小正特征值.

当 $\lambda$ 穿过 $A$ 的谱值 $\lambda_i$ 时, 指数 $i(I-\lambda A)$ 乘以 $(-1)^m$, 其中 $m$ 是 $\lambda_i$ 的次数, 即 $\operatorname{Ker}(I-\lambda_iA)^k$ 的维数; 当 $k$ 充分大时, 该维数与 $k$ 无关.

\subsection{非唯一性定理}
\label{subsec:ch2-taylor-nonuniqueness-theorem}
我们的目标是证明以下结果.

\MBSourceStatementNumber{theorem}{1}
\begin{Theorem}
\label{thm:ch2-taylor-nonuniqueness}
当 $\lambda$ 充分大且 $L$ 取适当值时, 问题 \eqref{eq:ch2-taylor-motion-equations}, \eqref{eq:ch2-taylor-incompressibility}, \eqref{eq:ch2-taylor-boundary-data} 有周期为 $L$ 的 $z$-周期解, 且该解不同于平凡解 \eqref{eq:ch2-taylor-trivial-base-solution}.
\end{Theorem}

\begin{Proof}
证明当 $\lambda$ 充分大时, 方程 \eqref{eq:ch2-taylor-functional-equation} 在 $V$ 中有非平凡解. 由引理 \ref{lem:ch2-taylor-simple-principal-eigenvalue}, 可以选择 $L$ 使 $\lambda_1$ 是 $B$ 在 $V$ 中的简单特征值; 这些就是定理 \ref{thm:ch2-taylor-nonuniqueness} 中所述的 $L$ 值.

由命题 \ref{prop:ch2-taylor-small-parameter-uniqueness} 可知, 当 $\lambda$ 充分小时($\lambda\leq c(r_1,r_2)$, 见 \eqref{eq:ch2-taylor-small-parameter-condition}), \eqref{eq:ch2-taylor-functional-equation} 只有平凡解.

若对某个 $\lambda\in[0,\lambda_1]$, 方程 \eqref{eq:ch2-taylor-functional-equation} 有非平凡解, 则定理已经得证. 因此以下假设
\begin{equation}
\phi=0\text{ 是任意 }\lambda\in[0,\lambda_1]\text{ 时 }\phi=\lambda T\phi\text{ 的唯一解}.
\label{eq:ch2-taylor-trivial-solution-assumption}
\end{equation}
在此假设下, 接下来的引理利用度理论证明存在 $V$ 中的非零 $\phi$, 满足 $\phi=\lambda T\phi$ 且 $\lambda>\lambda_1$.
\end{Proof}

\MBSourceStatementNumber{lemma}{14}
\begin{Lemma}
\label{lem:ch2-taylor-degree-boundary-local}
设 $\omega$ 是以 $0$ 为中心的 $V$ 的某个开球. 存在 $\delta>0$, 使对区间 $[\lambda_1,\lambda_1+\delta]$ 中每个 $\lambda$, 方程 $\phi=\lambda T\phi$ 在 $\omega$ 的边界 $\partial\omega$ 上无解.
\end{Lemma}

\begin{Proof}
反证. 若结论不成立, 则存在递减到 $\lambda_1$ 的序列 $\lambda_n$ 和属于 $\partial\omega$ 的序列 $u_n$, 使 $u_n=\lambda_nTu_n$. 因为 $u_n$ 有界, 由引理 \ref{lem:ch2-taylor-operator-compact}, $Tu_n$ 相对紧, 存在子序列 $Tu_{n_i}$ 收敛到 $V$ 中某个极限 $v$. 于是 $u_{n_i}=\lambda_{n_i}Tu_{n_i}$ 收敛到 $\lambda_1v$. 由 $T$ 的连续性,
\[
\lambda_1v=\lambda_1T(\lambda_1v).
\]
所以 $\lambda_1v$ 是 $\phi=\lambda_1T\phi$ 的解, 由 \eqref{eq:ch2-taylor-trivial-solution-assumption}, $\lambda_1v=0$, $v=0$. 这与 $\lVert\lambda_1v\rVert_V$ 等于球 $\omega$ 的半径相矛盾(对所有 $n$, $\lVert u_n\rVert$ 等于 $\omega$ 的半径).
\end{Proof}

\MBSourceStatementNumber{lemma}{15}
\begin{Lemma}
\label{lem:ch2-taylor-degree-boundary-global}
在假设 \eqref{eq:ch2-taylor-trivial-solution-assumption} 下, 若 $\omega,\delta$ 如引理 \ref{lem:ch2-taylor-degree-boundary-local} 所述, 则对任意 $\lambda\in[0,\lambda_1+\delta)$, 方程 $\phi=\lambda T\phi$ 在 $\partial\omega$ 上无解.
\end{Lemma}

这是引理 \ref{lem:ch2-taylor-degree-boundary-local} 的显然推论. 该引理使我们可以对 $\lambda\in[0,\lambda_1+\delta)$ 定义度 $d(I-\lambda T,\omega,0)$.

\hypertarget{chapter-02-page-180}{}
\MBSourceStatementNumber{lemma}{16}
\begin{Lemma}
\label{lem:ch2-taylor-degree-one}
在上述 $\delta,\omega$ 下,
\[
d(I-\lambda T,\omega,0)=1,
\qquad \lambda\in[0,\lambda_1+\delta).
\]
\end{Lemma}

\begin{Proof}
由拓扑度的性质 (iii),
\[
d(I-\lambda T,\omega,0)=d(I,\omega,0)=i(I),
\]
而该指数等于 $1$(恒等算子的指数).
\end{Proof}

\MBSourceStatementNumber{lemma}{17}
\begin{Lemma}
\label{lem:ch2-taylor-nontrivial-solution-near-bifurcation}
在假设 \eqref{eq:ch2-taylor-trivial-solution-assumption} 下, 对任意 $\lambda\in(\lambda_1,\lambda_1+\delta)$, 至少存在一个 $\phi=\lambda T\phi$ 的非平凡解.
\end{Lemma}

\begin{Proof}
由引理 \ref{lem:ch2-taylor-principal-degree-one} 和指数的性质 (iv), $i(I-\lambda B)$ 在 $[0,\lambda_1)$ 上等于 $1$, 在 $(\lambda_1,\lambda_1+\delta)$ 上等于 $-1$. 由指数的性质 (iii), $i(I-\lambda T,0)$ 在 $(0,\lambda_1)$ 上为 $+1$, 在 $(\lambda_1,\lambda_1+\delta)$ 上为 $-1$. 若 $\lambda\in(\lambda_1,\lambda_1+\delta)$ 且 $0$ 是 $\phi=\lambda T\phi$ 的唯一解, 则由指数的性质 (i) 应有
\[
d(I-\lambda T,\omega,0)=i(I-\lambda T,0).
\]
但我们已经证明
\[
d(I-\lambda T,\omega,0)=+1,
\qquad
i(I-\lambda T,0)=-1,
\qquad
\lambda\in(\lambda_1,\lambda_1+\delta).
\]
因此对任意 $\lambda\in(\lambda_1,\lambda_1+\delta)$, 方程 $\phi=\lambda T\phi$ 有非平凡解. 定理 \ref{thm:ch2-taylor-nonuniqueness} 的证明完成.
\end{Proof}

\MBSourceStatementNumber{remark}{3}
\begin{Remark}
\label{rem:ch2-taylor-large-lambda-physical-meaning}
条件"$\lambda$ 充分大"意味着角速度 $\alpha$ 大, 或者在 $r_1,r_2$ 固定时黏性 $\nu$ 小.
\end{Remark}

\MBSourceStatementNumber{remark}{4}
\begin{Remark}
\label{rem:ch2-taylor-bifurcation}
在条件 \eqref{eq:ch2-taylor-trivial-solution-assumption} 下, 对每个 $\lambda\in(\lambda_1,\lambda_1+\delta)$, \eqref{eq:ch2-taylor-functional-equation} 有一个非平凡解 $\phi_\lambda$. 可以证明, 当 $\lambda$ 递减到 $\lambda_1$ 时, $\phi_\lambda\to0$ 于 $V$. 这就是分岔.
\end{Remark}

对 Bénard 问题, 情形非常类似, 但可以证明当 $\lambda\in[0,\lambda_1]$ 时只有平凡解. 因此不需要假设 \eqref{eq:ch2-taylor-trivial-solution-assumption}, 并且确实证明了分岔的出现(见 \MBUnresolvedReference{citation}{V.I. Yudovich [2]}, \MBUnresolvedReference{citation}{Rabinowitz [1]}, \MBUnresolvedReference{citation}{Velte [1]}).

\MBUnresolvedReference{citation}{Rabinowitz [5]} 用分析方法研究了 Taylor 问题.

\paragraph{致谢.}
作者衷心感谢 P.H. Rabinowitz 对本节提出的有益意见.

\hypertarget{chapter-02-page-181}{}
